
the \(\mathscr C\)-group structure on \(G\), and finally by
\(\operatorname{pr}_{2,3}\) the projection of
\(G\times G\times X=G\times(G\times X)\) onto the second factor. Then

\SGAThreeDiagram[0.47\textwidth]
  {figures/exp5/p256_action_groupoid_continuation.png}
  {The groupoid attached to the action of the \(\mathscr C\)-group \(G\) on \(X\).}
  {fig:sga3-exp5-p256-action-groupoid}
  {LNM 151, p.~256}

\noindent is a \(\mathscr C\)-groupoid.

\medskip
\phantomsection\label{par:sga3-V-2-b}%
\noindent\textbf{b)} Let \(d_0,d_1:X_1\to X_0\) be an
\emph{equivalence pair}. In other words, if
\(d_0\boxtimes d_1:X_1\to X_0\times X_0\) is the arrow with components
\(d_0\) and \(d_1\), we suppose that, for every object \(T\) of
\(\mathscr C\), the map \((d_0\boxtimes d_1)(T)\) is a bijection from
\(X_1(T)\) onto the graph of an equivalence relation on \(X_0(T)\).
Consequently, \(X_1(T)\) is identified with the set of pairs \((x,y)\) of
elements of \(X_0(T)\) such that \(x\sim y\). Likewise,
\[
  X_2(T)=\bigl(X_1\times_{d_1,d_0}X_1\bigr)(T)
\]
is identified with the set of triples \((x,y,z)\) of elements of
\(X_0(T)\) such that \(x\sim y\) and \(y\sim z\). There is therefore a
unique arrow \(d'_1:X_2\to X_1\) making squares
\hyperref[sga3:v:lemma:1-2:equation:1]{(1)} and
\hyperref[sga3:v:lemma:1-2:equation:2]{(2)} commute: \(d'_1(T)\) must send
\((x,y,z)\in X_2(T)\) to \((x,z)\in X_1(T)\). With this choice of
\(d'_1\),
\[
  \bigl(d_0,d_1:X_1\rightrightarrows X_0,d'_1\bigr)
\]
is a \(\mathscr C\)-groupoid.

Conversely, consider a \(\mathscr C\)-groupoid \(X_\ast\) for which
\(d_0\boxtimes d_1:X_1\to X_0\times X_0\) is a monomorphism. Then
\((d_0,d_1)\) is an equivalence pair, and \(X_\ast\) can be reconstructed
from \((d_0,d_1)\) as explained \hyperref[par:sga3-V-2-b]{a few lines above}.%

\medskip
\phantomsection\label{par:sga3-V-2-c}%
\noindent\textbf{c)} If \(p:X\to Y\) is any arrow of \(\mathscr C\), and
if \(\operatorname{pr}_1\) and \(\operatorname{pr}_2\) are the two
projections of \(X\times_{p,p}X\) onto \(X\), then
\[
  (\operatorname{pr}_1,\operatorname{pr}_2):
  X\times_{p,p}X\rightrightarrows X
\]
is an equivalence pair. The arrow \(p\) is called an
\emph{effective epimorphism} if the diagram

\SGAThreeDiagram[0.36\textwidth]
  {figures/exp5/p256_effective_epimorphism.png}
  {The defining diagram for an effective epimorphism.}
  {fig:sga3-exp5-p256-effective-epimorphism}
  {LNM 151, p.~256}

\noindent is exact, that is, if
\[
  (Y,p)=\operatorname{Coker}(\operatorname{pr}_1,\operatorname{pr}_2).
\]

For example, let \(S\) be a Noetherian scheme and let \(\mathscr C\) be
the category of schemes finite over \(S\). We show that an epimorphism in
\(\mathscr C\) need not be effective. Take
\[
  S=\operatorname{Spec}k[T^3,T^5],\qquad
  Y=S,\qquad X=\operatorname{Spec}k[T],
\]
where \(k\) is a commutative field. If \(i\) is the inclusion of
\(B=k[T^3,T^5]\) into \(A=k[T]\), take \(p=\operatorname{Spec}i\). Then
\(X\times_{p,p}X\) is identified with \(\operatorname{Spec}(A\otimes_B A)\),
and \(\operatorname{Coker}(\operatorname{pr}_1,\operatorname{pr}_2)\)
with \(\operatorname{Spec}B'\), where \(B'\) is the subring of \(A\)
consisting of the elements \(a\) such that
\(a\otimes_B1=1\otimes_Ba\). But
\[
  T^7\otimes_B1
  =(T^2T^5)\otimes_B1
  =T^2\otimes_BT^5
  =T^2\otimes_B(T^3T^2)
  =T^5\otimes_BT^2
  =1\otimes_BT^7.
\]
Thus \(T^7\) belongs to \(B'\), does not belong to \(B\), and
\(\operatorname{Spec}B'\) is different from \(\operatorname{Spec}B\),
which gives the counterexample.%

\section*{3. Some sorites concerning \(\mathscr C\)-groupoids}
\phantomsection\label{sec:sga3-V-3}
\addcontentsline{toc}{section}{3. Some sorites concerning
  \texorpdfstring{\(\mathscr C\)}{C}-groupoids}

Here are, in no particular order, some remarks used below.

\medskip
\phantomsection\label{par:sga3-V-3-a}%
\noindent\textbf{a)} Let

\SGAThreeDiagram[0.31\textwidth]
  {figures/exp5/p257_x_groupoid.png}
  {The \(\mathscr C\)-groupoid \(X_\ast\).}
  {fig:sga3-exp5-p257-x-groupoid}
  {LNM 151, pp.~257--258}

\noindent be a \(\mathscr C\)-groupoid, and let \(f_0:Y_0\to X_0\) be
an arrow of \(\mathscr C\). We shall define a \(\mathscr C\)-groupoid
with base \(Y_0\),

\SGAThreeDiagram[0.31\textwidth]
  {figures/exp5/p257_y_induced_groupoid.png}
  {The induced \(\mathscr C\)-groupoid \(Y_\ast\).}
  {fig:sga3-exp5-p257-y-induced-groupoid}
  {LNM 151, pp.~257--258}

\noindent which will be said to be \emph{induced by \(X_\ast\) and
\(f_0\)}. We also say that \(Y_\ast\) is the \emph{inverse image of
\(X_\ast\) under the base change \(f_0\)}.

We choose \(Y_1\) to be the fiber product of the diagram

\SGAThreeDiagram[0.29\textwidth]
  {figures/exp5/p257_base_change_fiber_square.png}
  {The fiber-product definition of \(Y_1\).}
  {fig:sga3-exp5-p257-base-change-fiber-square}
  {LNM 151, pp.~257--258}

\noindent and take \(e_0\) and \(e_1\) to be the composites of the
canonical arrow \(Y_1\to Y_0\times Y_0\) with the first and second
projections of \(Y_0\times Y_0\), respectively. The morphism
\(Y_1\to Y_0\times Y_0\) is then \(e_0\boxtimes e_1\), and
\[
  f_0\circ e_i=d_i\circ f_1\qquad(i=0,1),
\]
where \(f_1\) denotes the projection of \(Y_1\) onto \(X_1\).

Set \(Y_2=Y_1\times_{e_0,e_1}Y_1\), cf.
\hyperref[par:sga3-V-1-c]{1(c)}. One may say that the pair \((e_0,e_1)\)
is defined so that, for every \(T\in\mathscr C\) and every pair \((y,x)\)
of elements of \(Y_0(T)\), there is a certain bijective correspondence
\[
  \psi\longmapsto{}_{y}\psi_{x}
\]
between the arrows \(\psi\) of \(X_\ast(T)\) with source \(f_0(x)\) and
target \(f_0(y)\), and the arrows \({}_{y}\psi_x\) of \(Y_\ast(T)\)
with source \(x\) and target \(y\). Thus \(e'_1:Y_2\to Y_1\) is
determined by defining, for every \(T\in\mathscr C\), the composition
of arrows in \(Y_\ast(T)\) by the formula
\[
  {}_{z}\psi_y\circ{}_{y}\varphi_x
  ={}_{z}(\psi\circ\varphi)_x.
\]
It is clear that this definition makes each \(Y_\ast(T)\) a groupoid.

\medskip
\phantomsection\label{par:sga3-V-3-b}%
\noindent\textbf{b)} Given the \(\mathscr C\)-groupoid \(X_\ast\) and
the base change \(f_0:Y_0\to X_0\), one can reconstruct the pair
\((e_0,e_1):Y_1\rightrightarrows Y_0\) in another way.%
Construct \(Y_0\times_{X_0}X_1\), \(\operatorname{pr}_1\), and
\(\operatorname{pr}_2\) so that the square

\SGAThreeDiagram[0.27\textwidth]
  {figures/exp5/p257_reconstruction_square.png}
  {The Cartesian square defining \(Y_0\times_{X_0}X_1\).}
  {fig:sga3-exp5-p257-reconstruction-square}
  {LNM 151, pp.~257--258}

\noindent is Cartesian. A straightforward reduction to the set-theoretic
case then shows that the following square is Cartesian:

\SGAThreeDiagram[0.27\textwidth]
  {figures/exp5/p258_cartesian_reconstruction_square.png}
  {The second Cartesian square in the reconstruction of \((e_0,e_1)\).}
  {fig:sga3-exp5-p258-cartesian-reconstruction-square}
  {LNM 151, p.~259}

Here \(f_1\) denotes the canonical projection of
\[
  Y_1=(Y_0\times Y_0)\times_{(X_0\times X_0)}X_1
\]
onto \(X_1\).

\medskip
\phantomsection\label{par:sga3-V-3-c}%
\noindent\textbf{c)} We give two examples of inverse images of a
\(\mathscr C\)-groupoid. Take \(Y_0=X_1\) and \(f_0=d_0\). For every
object \(T\) of \(\mathscr C\), the set \(Y_1(T)\) is then identified
with the set of diagrams of the form

\SGAThreeDiagram[0.18\textwidth]
  {figures/exp5/p258_inverse_image_phi_diagram.png}
  {An arrow of the first inverse-image groupoid.}
  {fig:sga3-exp5-p258-inverse-image-phi}
  {LNM 151, p.~259}

\noindent in \(X_\ast(T)\). The source of such a diagram is the arrow
\(f\), and its target is the arrow \(g\). These diagrams compose in the
evident way.

Now set \(Y'_0=X_1\) and \(f'_0=d_1\) (we add the primes% 
to avoid any confusion with the preceding example). In this case, for
every \(T\in\mathscr C\), the set \(Y'_1(T)\) is identified with the set
of diagrams of the form

\SGAThreeDiagram[0.18\textwidth]
  {figures/exp5/p258_inverse_image_psi_diagram.png}
  {An arrow of the second inverse-image groupoid.}
  {fig:sga3-exp5-p258-inverse-image-psi}
  {LNM 151, p.~259}

\noindent in the groupoid \(X_\ast(T)\). The source of such a diagram is
\(f\), its target is \(g\), and these diagrams again compose in the
evident way.

It is now clear that the identity map of \(Y_0(T)\), together with the
map

\SGAThreeDiagram[0.44\textwidth]
  {figures/exp5/p258_groupoid_isomorphism_map.png}
  {The map from \(Y_1(T)\) to \(Y'_1(T)\).}
  {fig:sga3-exp5-p258-groupoid-isomorphism-map}
  {LNM 151, p.~259}

\noindent defines an isomorphism from the groupoid \(Y_\ast(T)\) onto
\(Y'_\ast(T)\). Moreover, this isomorphism depends functorially on \(T\),
so the \(\mathscr C\)-groupoids \(Y_\ast\) and \(Y'_\ast\) are
isomorphic.%

\medskip
\phantomsection\label{par:sga3-V-3-d}%
\noindent\textbf{d)}

\medskip
\phantomsection\label{prop:sga3-V-3.1}%
\noindent\textbf{Proposition 3.1.}---
\textit{Keep the notation of \hyperref[par:sga3-V-3-a]{3(a)}, and suppose that \(f_0\) is an
effective and universal epimorphism. Then
\(\operatorname{Coker}(d_0,d_1)\) exists if and only if
\(\operatorname{Coker}(e_0,e_1)\) exists.%
\footnote{\label{note:sga3-V-14}Editor's note: the original has been
modified to display explicitly the isomorphism below.}
Moreover, in this case \(f_0\) induces an isomorphism}
\phantomsection\label{eq:sga3-V-3.1-cokernel-isomorphism}%
\[
  \operatorname{Coker}(d_0,d_1)
  \xrightarrow{\ \sim\ }
  \operatorname{Coker}(e_0,e_1).
\]

Recall first that an epimorphism \(f_0:Y_0\to X_0\) is called
\emph{universal} if, in every Cartesian square

\SGAThreeDiagram[0.18\textwidth]
  {figures/exp5/p259_universal_epimorphism_square.png}
  {A base change of the universal epimorphism \(f_0\).}
  {fig:sga3-exp5-p259-universal-epimorphism-square}
  {LNM 151, p.~260}

\noindent the arrow \(f'\) is an epimorphism. Now let \(C(d_0,d_1)\)
denote the covariant functor from \(\mathscr C\) to sets which assigns to
every \(T\in\mathscr C\) the kernel of the pair
\[
  T(d_0),T(d_1):T(X_0)\rightrightarrows T(X_1),
\]
and define \(C(e_0,e_1)\) in the same way. For every
\(T\in\mathscr C\), there is therefore a commutative diagram

\SGAThreeDiagram[0.51\textwidth]
  {figures/exp5/p259_functor_kernel_diagram.png}
  {The comparison of the functors \(C(d_0,d_1)\) and \(C(e_0,e_1)\).}
  {fig:sga3-exp5-p259-functor-kernel-diagram}
  {LNM 151, p.~260}

Here \(T(f)\) is the injection induced by the injection \(T(f_0)\). If
we show that \(T(f)\) is surjective for every \(T\), we obtain a
functorial isomorphism
\[
  f:C(d_0,d_1)\xrightarrow{\ \sim\ }C(e_0,e_1).
\]
Thus one of these functors is representable if and only if the other is,
which proves \hyperref[prop:sga3-V-3.1]{the proposition}.

To prove that \(T(f)\) is surjective, consider the diagram

\SGAThreeDiagram[0.52\textwidth]
  {figures/exp5/p259_delta_proof_diagram.png}
  {The diagram used to prove the surjectivity of \(T(f)\).}
  {fig:sga3-exp5-p259-delta-proof-diagram}
  {LNM 151, p.~260}

Here \(\Delta\) is the section of \(Y_1\to Y_0\times Y_0\) defined by
the morphism
\[
  s\circ f_0\circ\operatorname{pr}_1:Y_0\times Y_0\longrightarrow X_1,
\]
where \(s:X_0\to X_1\) satisfies
\hyperref[eq:sga3-exp5-3]{(3)} and
\hyperref[eq:sga3-exp5-3bis]{(3 bis)} of \hyperref[sec:sga3-V-1]{Section~1}.

If \(g:Y_0\to T\) satisfies \(g\circ e_0=g\circ e_1\), then
\[
  g\circ e_0\circ\Delta=g\circ e_1\circ\Delta,
\]
so \(g\circ\operatorname{pr}_1=g\circ\operatorname{pr}_2\). Since
\(f_0\) is an effective epimorphism, \(g\) factors as \(h\circ f_0\)
for an arrow \(h:X_0\to T\); in other words,
\(g=T(f_0)(h)\). It remains to prove that
\(h\in C(d_0,d_1)(T)\), or equivalently that \(hd_0=hd_1\). But
\[
  hd_0f_1=hf_0e_0=ge_0=ge_1=hf_0e_1=hd_1f_1,
\]
and the desired equality follows because \(f_1\) is an epimorphism
(since \(f_0\) is a universal epimorphism).

\medskip
\phantomsection\label{par:sga3-V-3-e}%
\noindent\textbf{e)} Now let \(S\) be a scheme and take
\(\mathscr C=(\mathbf{Sch}/S)\). The data of a \(\mathscr C\)-groupoid

\SGAThreeDiagram[0.37\textwidth]
  {figures/exp5/p260_scheme_groupoid.png}
  {A \((\mathbf{Sch}/S)\)-groupoid.}
  {fig:sga3-exp5-p260-scheme-groupoid}
  {LNM 151, pp.~261--262}

\noindent defines an equivalence relation on the set \(\underline{X_0}\)
underlying the scheme \(X_0\): for \(x,y\in\underline{X_0}\), write
\(x\sim y\) if there exists \(z\in\underline{X_1}\) such that
\(x=d_1z\) and \(y=d_0z\). Reflexivity and symmetry are clear.%
\footnote{\label{note:sga3-V-15}Editor's note: reflexivity follows from
the existence of \(s:X_0\to X_1\), which is a section of both \(d_0\)
and \(d_1\); symmetry follows from the existence of the involution
\(\sigma\) of \(X_1\) which ``exchanges \(d_0\) and \(d_1\)'', that is,
which satisfies \(d_0\sigma=d_1\) and \(d_1\sigma=d_0\); cf. \hyperref[sec:sga3-V-1]{Section~1},
\hyperref[eq:sga3-exp5-3]{(3)},
\hyperref[eq:sga3-exp5-3bis]{(3 bis)}, and
\hyperref[eq:sga3-exp5-dagger]{(\(\dagger\))}.}
Let us prove transitivity. If \(x\sim y\) and \(y\sim z\), there exist
\(u,v\in\underline{X_1}\) such that
\[
  x=d_1u,\qquad y=d_0u,\qquad y=d_1v,\qquad z=d_0v.
\]
It follows that \((v,u)\) belongs to the set-theoretic fiber product
\(\underline{X_1}\times_{\underline{d_1},\underline{d_0}}
\underline{X_1}\). Since the canonical map

\SGAThreeDiagram[0.32\textwidth]
  {figures/exp5/p260_underlying_fiber_product_map.png}
  {The canonical map from the underlying set of the fiber product to the
  fiber product of the underlying sets.}
  {fig:sga3-exp5-p260-underlying-fiber-product-map}
  {LNM 151, pp.~261--262}

\noindent is surjective, \((v,u)\) is the image of some
\(w\in\underline{X_2}\). Hence
\[
  x=d_1d'_1w,\qquad z=d_0d'_1w,
\]
and therefore \(x\sim z\).

\medskip
\phantomsection\label{par:sga3-V-3-f}%
\noindent\textbf{f)} Keep the notation of \hyperref[par:sga3-V-3-a]{3(a)} and \hyperref[par:sga3-V-3-b]{3(b)}, still taking
\(\mathscr C=(\mathbf{Sch}/S)\). If \(x,y\) are points of \(Y_0\), then
\[
  x\sim y\quad\Longleftrightarrow\quad f_0(x)\sim f_0(y).
\]
Thus the inverse image of the equivalence relation defined by a groupoid
is the equivalence relation defined by the inverse image of that
groupoid.

Indeed, suppose that \(x\sim y\). There is a \(z\in\underline{Y_1}\)
such that \(x=e_1(z)\) and \(y=e_0(z)\). Since
\(f_0\circ e_i=d_i\circ f_1\) for \(i=0,1\), it follows that
\[
  f_0(x)=d_1f_1(z),\qquad f_0(y)=d_0f_1(z),
\]
and hence \(f_0(x)\sim f_0(y)\).

Conversely, suppose that \(f_0(x)\sim f_0(y)\), and let
\(z\in\underline{X_1}\) satisfy
\[
  f_0(y)=d_1(z),\qquad f_0(x)=d_0(z).
\]
Using the construction and notation of \hyperref[par:sga3-V-3-b]{3(b)}, there is a point
\(t\in\underline{Y_0\times_{X_0}X_1}\) such that
\(\operatorname{pr}_1(t)=x\) and \(\operatorname{pr}_2(t)=z\). Since
\(f_0(y)=d_1\operatorname{pr}_2(t)\), there is likewise an
\(s\in\underline{Y_1}\) such that
\[
  y=e_1(s),\qquad (e_0\boxtimes f_1)(s)=t.
\]
It follows that
\[
  e_0(s)=\operatorname{pr}_1(e_0\boxtimes f_1)(s)
  =\operatorname{pr}_1(t)=x,
\]
so \(x\sim y\).
